% !TEX root = task-2-1-1.tex
% Editable task bank. Append new BankTask blocks at the end of the document.
\documentclass[12pt,a4paper]{article}
\usepackage{physics-textbook}
\usepackage{graphicx,needspace}
\usetikzlibrary{decorations.pathmorphing}
\newif\ifShowSolutions
\ifdefined\TutorEdition
  \ShowSolutionstrue
\else
  \ShowSolutionsfalse % Change to \ShowSolutionstrue to display solutions.
\fi
\newcounter{banktask}
\newsavebox{\BankTaskBox}
\newcommand{\BankTaskContent}[3]{%
  \noindent\textbf{\thebanktask.}\quad #2\par
  \ifShowSolutions
    \Needspace{50mm}% Keep the solution heading with its opening text or diagram.
    #3\par
    \ifdefined\TutorEdition
      \ifcsname TutorSource#1\endcsname
        \nopagebreak\smallskip{\small\csname TutorSource#1\endcsname\par}%
      \fi
    \fi
  \fi
}
\newcommand{\BankTask}[3]{%
  \refstepcounter{banktask}\par\medskip
  \sbox{\BankTaskBox}{\begin{minipage}{\linewidth}%
    \BankTaskContent{#1}{#2}{#3}\end{minipage}}%
  \ifdim\dimexpr\ht\BankTaskBox+\dp\BankTaskBox\relax>\dimexpr\textheight-20mm\relax
    \Needspace{35mm}\BankTaskContent{#1}{#2}{#3}%
  \else
    \Needspace{\dimexpr\ht\BankTaskBox+\dp\BankTaskBox\relax}%
    \noindent\usebox{\BankTaskBox}\par
  \fi
  \medskip
}
\begin{document}
\BookSubsection{2.1.1}{Атомы и молекулы}


% Source sdamgia-806: https://phys-ege.sdamgia.ru/problem?id=806
\BankTask{sdamgia-806}{
Из контейнера с твердым литием изъяли $4\,\text{моль}$ этого вещества. При этом число атомов лития в контейнере уменьшилось на

$1$) $4 \cdot 10 ^{23}$ 

$2$) $12 \cdot 10 ^{23}$ 

$3$) $24 \cdot 10 ^{23}$ 

$4$) $36 \cdot 10 ^{23}$
}{
\smallskip\noindent\textbf{Решение}\par
В одном моле любого вещества содержится $6 \cdot 10 ^{23}$ частиц. Таким образом, после изъятия $4\,\text{моль}$ вещества число атомов лития в контейнере уменьшилось на 

$4 \cdot 6 \cdot 10 ^{23} =24 \cdot 10 ^{23} .$

Ответ: $3$.
}

% Source sdamgia-3361: https://phys-ege.sdamgia.ru/problem?id=3361
\BankTask{sdamgia-3361}{
В баллоне находится $0{,}01\,\text{моль}$ газа. Сколько примерно молекул газа находится в баллоне?

$1$) $10 ^{21}$

$2$) $6 \cdot 10 ^{21}$

$3$) $10 ^{2}$

$4$) $6 \cdot 10 ^{24}$
}{
\smallskip\noindent\textbf{Решение}\par
В одном моль любого вещества находится приблизительно $N_A=6 \cdot 10 ^{23}$ молекул. Следовательно, в баллоне находится примерно $0{,}01 \cdot 6 \cdot 10 ^{23} =6 \cdot 10 ^{21}$ молекул.

Ответ: $2$.
}

% Source sdamgia-936: https://phys-ege.sdamgia.ru/problem?id=936
\BankTask{sdamgia-936}{
Плотность алюминия в $3$ раза больше плотности льда. В $1$ моле алюминия содержится:

$1$) атомов в $3$ раза больше, чем в одном моле льда;

$2$) столько же атомов, сколько в одном моле льда;

$3$) в $3$ раза меньше атомов, чем в одном моле льда;

$4$) на $12 \cdot 10 ^{23}$ атомов больше, чем в одном моле льда.
}{
\smallskip\noindent\textbf{Решение}\par
Количество вещества — это физическая величина, характеризующая количество однотипных структурных единиц, содержащихся в веществе. То есть эта величина «пересчитывает» количество молекул (для структурных веществ) и количество атомов (для «простых» веществ). 

В одном моле любого вещества содержится столько же структурных единиц вещества, сколько атомов содержится в $12$ граммах углерода. Следовательно, в одном моле алюминия содержится столько же атомов, сколько молекул в одном моле льда. Но молекула воды состоит из трех атомов, а значит, количество атомов в одном моле алюминия в $3$ раза меньше, чем количество атомов в одном моле льда.

Ответ: $3$.
}

% Source sdamgia-3769: https://phys-ege.sdamgia.ru/problem?id=3769
\BankTask{sdamgia-3769}{
В сосуде А находится $28\,\text{г}$ молекулярного азота, а в сосуде Б — $44\,\text{г}$ углекислого газа. В каком сосуде находится больше атомов?

$1$) в сосуде А

$2$) в сосуде Б

$3$) В сосудах А и Б содержится примерно одинаковое число атомов

$4$) в сосуде, объем которого больше
}{
\smallskip\noindent\textbf{Решение}\par
Это задание — удобный повод вспомнить, что атомы и молекулы — не одно и тоже. 

Определим, какое количество газов содержится в обоих сосудах. Поскольку молярная масса молекулярного азота равна $28\text{г}/\text{моль},$ а молекулярная масса углекислого газа — $44\text{г}/\text{моль},$ заключаем, что в обоих сосудах по одному моль вещества. Количество вещества — это физическая величина, характеризующая количество однотипных структурных единиц, содержащихся в веществе. То есть эта величина «пересчитывает» количество молекул. Если теперь вспомнить, что молекула азота $( N_2 )$ состоит из двух атомов, а молекула углекислого газа $( CO_2 )$ — из трех, то сразу ясно, что атомов в сосуде с углекислым газом приблизительно в полтора раза больше.

Ответ: $2$.
}


% Source sdamgia-3739: https://phys-ege.sdamgia.ru/problem?id=3739
\BankTask{sdamgia-3739}{
В двух сосудах находится по одному молю разных идеальных газов. Можно утверждать, что

$1$) число молекул, также как и число атомов в этих сосудах одинаково

$2$) число атомов в этих сосудах одинаково

$3$) число молекул в этих сосудах может быть различным

$4$) число атомов в этих сосудах может быть различным
}{
\smallskip\noindent\textbf{Решение}\par
Количество вещества — это физическая величина, характеризующая количество однотипных структурных единиц, содержащихся в веществе. То есть эта величина «пересчитывает» количество молекул. В $1$ моле любого газа всегда одинаковое количество молекул, равное числу Авогадро, а вот атомов может быть разное количество. Например, если один газ одноатомный, то в нем атомов столько же, сколько и молекул, а для двухатомного газа — количество атомов вдвое больше, чем количество молекул. Таким образом, среди приведенных утверждений всегда справедливо только утверждение под номером $4$.

Ответ: $4$.
}

% Source sdamgia-3209: https://phys-ege.sdamgia.ru/problem?id=3209
\BankTask{sdamgia-3209}{
Какая масса у $3\,\text{моль}$ неона?

$1$) $12\,\text{г}$

$2$) $36\,\text{г}$

$3$) $60\,\text{г}$

$4$) $18 \cdot 10 ^{- 23} \text{г}$
}{
\smallskip\noindent\textbf{Решение}\par
Молярная масса неона равна $\mu _{\mathrm{Ne}}=20 \cdot 10 ^{- 3} \frac{\text{кг}}{\text{моль}} .$Следовательно $3\,\text{моль}$ неона имеют массу $3\text{моль}\!\!\!\! \cdot 20 \cdot 10 ^{- 3} \frac{\text{кг}}{\text{моль}} =60\text{г}.$

Ответ: $3$.
}

% Source sdamgia-19661: https://phys-ege.sdamgia.ru/problem?id=19661
\BankTask{sdamgia-19661}{
Порция идеального одноатомного газа при температуре $300\,\text{К}$ обладает внутренней энергией $1{,}242\,\text{Дж}$. Сколько атомов содержит эта порция газа? (В $10^{20}$; целое число.)
}{
\smallskip\noindent\textbf{Решение}\par
Внутренняя энергия идеального одноатомного газа: 

$U = \frac{3}{2} \nu RT = \frac{3}{2} \cdot \frac{N}{N_A} RT.$

Откуда и находим число атомов: 

$N = \frac{2UN_A}{3RT} = 2 \cdot 10 ^{20} .$

Ответ: $2$.
}

% Source sdamgia-43492: https://phys-ege.sdamgia.ru/problem?id=43492
\BankTask{sdamgia-43492}{
Пористая неподвижная перегородка расположена посередине сосуда. В его левой половине находится $4\,\text{моль}$ гелия, а в правой — $40\,\text{г}$ аргона. Перегородка является непроницаемой для молекул аргона и может свободно пропускать лишь молекулы гелия. Температуры газов равны и постоянны. Из приведенного ниже списка

$1$) Концентрация гелия в правой половине сосуда в $2$ раза меньше концентрации аргона.

$2$) Давление газов в правой половине сосуда в $1{,}5$ раза больше давления газа в левой половине.

$3$) В правой половине сосуда общее число молекул газов больше, чем в левой половине.

$4$) Среднеквадратичная скорость теплового движения молекул аргона больше, чем у молекул гелия.

$5$) В результате установления равновесия давление газа в левой половине сосуда увеличилось в $2$ раза.
}{
\smallskip\noindent\textbf{Решение}\par
Найдем количество аргона: $\nu= \frac{m}{M} = \frac{40 \cdot 10 ^{- 3}}{40 \cdot 10 ^{- 3}} =1\text{моль}.$

$1$) Неверно. Концентрация газа $n= \frac{N}{V} = \frac{\nu N_A}{V} .$ Так как количество вещества гелия больше количества вещества аргона в $4$ раза, то и концентрация гелия больше концентрации аргона в $4$ раза.

$2$) Верно. Ввиду проницаемости перегородки в левой части будет содержаться $2$ моля вещества, а в правой — $3$ моля вещества. Из уравнения Клапейрона-Менделеева $pV=\nu RT$ следует, что давление газов в правой половине в $1{,}5$ раза больше, чем в левой.

$3$) Верно. В левой части будет содержаться $2$ моля вещества, а в правой — $3$ моля вещества, значит, число молекул в правой половине больше, чем в левой.

$4$) Неверно. Из формулы $\overline v = \sqrt{\frac{RT}{M}}$ следует, что среднеквадратичная скорость движения молекул аргона меньше, чем у молекул гелия.

$5$) Неверно. Количество вещества в левой половине уменьшилось в $2$ раза, следовательно, и давление уменьшилось в $2$ раза.

Ответ: $23$.
}

% Source sdamgia-30494: https://phys-ege.sdamgia.ru/problem?id=30494
\BankTask{sdamgia-30494}{
В сосуде объемом $V$ находится идеальный одноатомный газ массой $m$ при температуре $T$. ($\mu$ — молярная масса газа, $R$ — универсальная газовая постоянная, $N_{A}$ — число Авогадро) 

\TaskMatchingColumns{Физическая величина}{$A$)  концентрация газа\par Б)  давление газа}{Формула}{$1$)  $\frac{mRT}{\mu V}$\par $2$)  $\frac{\mu RT}{m V}$\par $3$)  $\frac{mN_A}{\mu V}$\par $4$)  $\frac{\mu N_A}{V}$}

\begin{center}\begin{tabular}{|c|c|}\hline
$A$ & Б\\\hline
\strut & \strut\\\hline\end{tabular}\end{center}
}{
\smallskip\noindent\textbf{Решение}\par
А) Концентрация по определению $n= \frac{N}{V} ,$ число молекул $N=\nu N_A= \frac{m}{\mu} N_A,$ тогда получаем $n= \frac{mN_A}{\mu V}$ (пункт $3$).

Б) Из уравнения Клапейрона — Менделеева $pV= \frac{m}{\mu} RT$ выражаем давление газа $p= \frac{mRT}{\mu V}$ (пункт $1$).

Ответ: $31$.
}

% Source sdamgia-4099: https://phys-ege.sdamgia.ru/problem?id=4099
\BankTask{sdamgia-4099}{
В сосуде объемом $1\,\text{л}$ находится одноатомный идеальный газ при давлении $2\,\text{кПа}$. Средняя кинетическая энергия теплового движения молекулы газа равна $6\cdot10^{-21}$ Дж. Сколько молекул газа содержится в этом сосуде? 

$1$) $1{,}55 \cdot 10 ^{22}$

$2$) $6{,}02 \cdot 10 ^{23}$

$3$) $0{,}33 \cdot 10 ^{21}$

$4$) $0{,}5 \cdot 10 ^{21}$
}{
\smallskip\noindent\textbf{Решение}\par
Идеальный газ подчиняется уравнению состояния Клапейрона-Менделеева: $pV=\nu RT.$ Средняя кинетическая энергия теплового движения молекул газа связан с температурой соотношением: $\overline{E}= \frac{3}{2} kT.$ Решая систему из этих двух уравнений и вспоминая связь между газовой постоянной, постоянной Больцмана и постоянной Авагадро $N_ak=R,$ для количества вещества в сосуде получаем: $\nu = \frac{3pVk}{2\overline{E}R} = \frac{3pV}{2\overline{E}N_a} .$ Отсюда, количество молекул газа в сосуде равно $N=\nu N_a= \frac{3pV}{2\overline{E}} = \frac{3 \cdot 2 \cdot 10 ^{3} \text{Па} \cdot 10 ^{- 3} \text{м} ^{3}}{2 \cdot 6 \cdot 10 ^{- 21} \text{Дж}} =0{,}5 \cdot 10 ^{21} .$

Ответ: $4$.
}

\end{document}
